\subsection{Гибридный выбор числа вызовов}\label{subsec:route-pilot}
Для RQ4 провели отдельную серию с~локальной Llama~3.2~3B Instruct (квантование Q4\_K\_M, Ollama 0.40.16, NVIDIA GeForce GTX~1060 с~6~ГБ памяти). Для неё составили 12 новых коротких задач на целочисленные функции Python. Например, требовалось вычислить смещение страницы, определить длину пересечения интервалов или округлить длительность в~секундах до целых минут в~сторону нуля. В~архивных карточках 1\,000 задач BigCodeBench отсутствовало поле с~требованием к~сессии исполнителя. В~новой серии это требование задали явно.

До генерации для шести задач в~протоколе задали требование отдельной сессии рецензента. Для остальных шести разрешили объединить этапы. Это организационное условие считывается из карточки задачи. Оно задаёт состав нагрузки пилота и~не оценивает сложность задачи. Требования параллельного выполнения в~наборе отсутствовали. Гибрид выбирал один или три вызова. Два контрольных способа всегда делали один или всегда три. Все три способа использовали одни и~те же названия ролей. Ещё в~двух условиях названия ролей заменили номерами этапов, сохранив формулировки действий и~число вызовов.

Каждая задача прошла все пять условий по одному разу. Порядок задач перемешали. Порядок условий внутри задачи сдвигали на одну позицию от задачи к~задаче, начиная со~случайной перестановки. Начальное состояние генератора зафиксировали, температуру задали равной нулю. В~отдельной сессии модель получала условие и~явно переданные ответы предшествующих этапов. На отдельный план отводилось до 40 выходных токенов, на черновой или итоговый код --- до~112. Гибрид для каждого назначения заново обращался к~модели. Всего выполнено 60 прогонов и~120 измеренных обращений после одного прогрева.

Во всех условиях модель должна была вернуть функцию Python с~одним оператором \texttt{return}. В~инструкции перечислялись допустимые операции над числами. Независимый оценщик проверял форму ответа и~совпадение результата с~ожидаемым целым числом на всех входах задачи. Всего в~12 наборах было 648 контрольных входов, включая граничные случаи. Тесты оставались скрытыми от модели. Время измеряли от выбора способа до окончания проверки, включая промежуточные запросы. Все обращения выполнялись последовательно на прогретой модели. Кэш между условиями сохранялся.

\begin{table}[!ht]
\caption{Локальный пилот. Показатели рассчитаны по одним и~тем же 12 задачам в~каждом условии}\label{tab:pilot}
\centering\small\renewcommand{\arraystretch}{1.22}
\begin{tabular}{@{}lrrrrr@{}}
\toprule
Способ & Успехи & Вызовы & Токены & Долл./1\,000 & Время, с \\
\midrule
Один с~ролями (SR) & 5/12 & 12 & 186,8 & 0,0612 & 0,66 \\
Гибрид (RouteMode) & 5/12 & 24 & 384,9 & 0,1272 & 1,35 \\
Три с~ролями (MAS) & 5/12 & 36 & 572,9 & 0,1876 & 1,93 \\
\midrule
Один с~номерами (SN) & 4/12 & 12 & 192,9 & 0,0636 & 0,68 \\
Три с~номерами (MN) & 4/12 & 36 & 608,2 & 0,2048 & 2,20 \\
\botrule
\end{tabular}
\end{table}
\FloatBarrier

Столбец <<Вызовы>> содержит общее число обращений за 12 прогонов. Токены (вход и~выход по счётчикам модели) и~время усреднены на один прогон, включая неуспешные. Для сопоставления цены взяты одинаковые условные ставки из расчёта без скидки на кэш, 0,20~доллара за миллион входных токенов и~1,20~доллара за миллион выходных. Стоимость дана в~пересчёте на 1\,000 прогонов. Генерация выполнена локально, денежные суммы представляют тарифный сценарий.

\textbf{Гибрид и~оба фиксированных способа с~ролями прошли тесты на одних и~тех же пяти задачах.} По сравнению с~постоянными тремя вызовами гибрид использовал в~среднем на 188 токенов меньше на задачу и~завершался на 0,58~с быстрее. По сравнению с~постоянным одним вызовом он израсходовал на 198 токенов больше и~занял на 0,70~с больше. Эти величины относятся к~нагрузке, в~которой отдельная сессия требовалась ровно в~половине случаев. Совпадение исходов на 12 задачах с~единственной генерацией на условие не устанавливает равенства качества на других задачах. Абсолютное время зависит от оборудования и~состояния кэша.

В~дополнительном сравнении названий ролей с~номерами этапов результат изменился на одной задаче из 12 в~обоих сравнениях. С~названиями ролей программа корректно вычислила номер позиции на кольце. При трёх нейтральных вызовах ответ дал неверное значение при положительном числе позиций в~кольце. При одном нейтральном вызове программа совпала с~эталоном на всех входах, но использовала два оператора \texttt{return}. Дополнительная проверка с~разрешённым ранним возвратом повысила результат SN с~4/12 до~5/12, показатели остальных условий сохранились. Таким образом, для одного вызова наблюдаемое различие относилось к~соблюдению формата ответа. На остальных 11 задачах исходы совпали. В~серии BigCodeBench соответствующие эффекты названий ролей составили $-0{,}75$ и~$+1{,}35$ процентного пункта (таблица~\ref{tab:scale-contrasts}). Это разные оценки для разных моделей и~наборов задач. Данные, полные запросы и~программа пересчёта приведены в~приложении~\ref{app:materials}.
